\documentclass[10pt,a4paper]{amsart}
\usepackage[margin=19mm]{geometry}
\usepackage{amsmath,amssymb,mathtools}
\usepackage[scr=rsfs,cal=euler]{mathalfa}
\usepackage{xfrac}
\usepackage[unicode,hidelinks,bookmarks=false]{hyperref}
\newcommand{\ie}{i.e.\ }
\newcommand{\cf}{cf.\ }
\newcommand{\resp}{respectively\ }
\newcommand{\Subsection}{\subsection}
\providecommand{\nobreakdash}{\nobreak\hbox{-}}
\setlength{\emergencystretch}{2em}
\multlinegap0pt
\usepackage{amssymb}
\usepackage{tikz}
\newtheorem{theorem}{Theorem}
\title{Nuclear dimension, pure infiniteness and real rank for higher rank graph $C^*$-algebras}
\author{David Pask}
\date{Source: arXiv:2607.27691v1 (CC BY 4.0)}
\begin{document}
\maketitle

\noindent\emph{Source and licence.} Nuclear dimension, pure infiniteness and real rank for higher rank graph $C^*$-algebras. Version arXiv:2607.27691v1 (CC BY 4.0), distributed by arXiv under the Creative Commons Attribution 4.0 International licence, http://creativecommons.org/licenses/by/4.0/, as recorded in the arXiv licence metadata for this exact version and shown on the abstract page https://arxiv.org/abs/2607.27691v1. Full licence notice: \texttt{licenses/arxiv-2607.27691-CC-BY-4.0.txt}.\par\smallskip

\noindent\textit{Editorial scope.} Counted unit: the article's theorem thm:approx (source0.tex lines 1217--1252). The article's complete original proof of it is delivered with it (source0.tex lines 1254--1342, one \texttt{proof} environment of 89 lines). No mathematics is written, repaired or re-derived: the statement and the proof are the article's own lines, in the article's own notation, and no retained line is substituted or re-flowed.  No further source-local context is needed: every cross-reference of every retained line resolves inside the two delivered spans.  Nothing was omitted from any retained span, so the package carries no omission record.  No retained line contains a citation, so the wrapper carries no bibliography and no bibliography entry.  No retained line contains a figure environment, an image-inclusion command or drawing code, so no image is delivered and the package declares no asset dependency.  Editorial formatting only, in addition: an AMS \texttt{amsart} preamble that restates the article's own declarations and macros used by the retained lines, namely the environments \texttt{theorem} on separate counters, together with the standard packages those lines need.  The retained environments print in this document as Theorem 1 in order of appearance, whereas the article prints the same items with the numbering of its own full document.  The complete native source of the article as distributed in the arXiv e-print for this exact version is bundled unchanged as \texttt{sources/206365/source0.tex}.
\begin{theorem}\label{thm:approx}
Let $\Lambda$ be a locally convex, row-finite $k$-graph with no sources.
Assume:
\begin{enumerate}
\item $\Lambda$ is strongly aperiodic;
\item every vertex in each maximal tail $M$ is connected to by a
      generalised cycle in $M$;
\item the lattice $H(\Lambda)$ of saturated hereditary subsets of
      $\Lambda^0$ is finite.
\end{enumerate}
Then there exists an increasing sequence of $C^*$-subalgebras


\[
A_1 \subset A_2 \subset \cdots \subset A_n = C^*(\Lambda)
\]


and a sequence of locally convex, row-finite, strongly aperiodic
$k$-graphs


\[
\Gamma_1 \subset \Gamma_2 \subset \cdots \subset \Gamma_n = \Lambda
\]


such that each $A_i \cong C^*(\Gamma_i)$ and


\[
\overline{\bigcup_{i=1}^n A_i} = C^*(\Lambda).
\]


\end{theorem}

\begin{proof}
Since $H(\Lambda)$ is finite, we may list its elements in increasing
order:


\[
\emptyset = H_0 \subsetneq H_1 \subsetneq \cdots \subsetneq H_m = \Lambda^0.
\]


For each $i$ let


\[
\Gamma_i := \Lambda\setminus H_{m-i},
\]


the full subgraph on the vertex set $\Lambda^0\setminus H_{m-i}$ with
all morphisms whose source lies in $\Lambda^0\setminus H_{m-i}$.
Because $H_{m-i}$ is saturated hereditary, $\Gamma_i$ is a locally
convex, row-finite $k$-graph with no sources.

\smallskip

\textbf{Claim 1: Each $\Gamma_i$ is strongly aperiodic.}
Since $\Lambda$ is strongly aperiodic, every quotient
$\Gamma(\Lambda\setminus H)$ is aperiodic for every saturated hereditary
$H\subset\Lambda^0$. But $\Gamma_i = \Gamma(\Lambda\setminus H_{m-i})$,
so each $\Gamma_i$ is aperiodic. Moreover, if $K\subset\Gamma_i^0$
is saturated hereditary in $\Gamma_i$, then $K\cup H_{m-i}$ is
saturated hereditary in $\Lambda$, and strong aperiodicity of
$\Lambda$ implies aperiodicity of $\Gamma(\Lambda\setminus (K\cup H_{m-i}))$.
Thus $\Gamma_i$ is strongly aperiodic.

\smallskip

\textbf{Claim 2: Each vertex of each maximal tail of $\Gamma_i$ is
connected to by a generalised cycle in $\Gamma_i$.}
Let $M$ be a maximal tail in $\Gamma_i$. Then $M$ is also a maximal
tail in $\Lambda$ (because adding vertices in $H_{m-i}$ cannot create
new paths into $M$). By hypothesis, every vertex of $M$ is connected
to by a generalised cycle in $\Lambda$. Since $M\subset\Gamma_i^0$,
and generalised cycles based at vertices of $M$ remain inside
$\Gamma_i$, the claim follows.

\smallskip

\textbf{Define the subalgebras.}
For each $i$ let


\[
A_i := C^*(\Gamma_i) \subset C^*(\Lambda),
\]


via the canonical inclusion of the Cuntz--Krieger family
$\{s_\lambda : \lambda\in\Gamma_i\}$ into $C^*(\Lambda)$.

\smallskip

\textbf{Claim 3: $A_i \subset A_{i+1}$ and $A_m = C^*(\Lambda)$.}
Since $\Gamma_i^0 \subset \Gamma_{i+1}^0$ and
$\Gamma_i^1 \subset \Gamma_{i+1}^1$, the universal property of
$C^*(\Gamma_i)$ gives an inclusion $A_i\subset A_{i+1}$.  
For $i=m$ we have $\Gamma_m = \Lambda$, so $A_m=C^*(\Lambda)$.

\smallskip

\textbf{Claim 4: $\overline{\bigcup_i A_i} = C^*(\Lambda)$.}
Every generator $s_\lambda$ of $C^*(\Lambda)$ lies in some
$A_i$, namely the smallest $i$ such that $s(\lambda)\notin H_{m-i}$.
Thus the Cuntz--Krieger family for $\Lambda$ is contained in
$\bigcup_i A_i$, and hence


\[
C^*(\Lambda) = \overline{\operatorname{span}}
\{ s_\lambda s_\mu^* : \lambda,\mu\in\Lambda\}
\subset \overline{\bigcup_i A_i}.
\]



\smallskip

Combining the claims proves the theorem.
\end{proof}

\end{document}
